\introduction\label{ch:intro}
        Nell'era digitale la figura del Tecnico Informatico è diventata oramai onnipresente ed è spesso fondamentale per il corretto funzionamento di aziende e di istituzioni. Le università in particolare necessitano della figura del Tecnico Informatico in quanto molti servizi offerti, se non tutti, prevedono componenti trasmesse per via informatica. \newline
        Il Tecnico Informatico è appunto quella figura che permette il corretto funzionamento dell'intera rete locale dell'Università e che sostiene il Corpo Docente e il personale amministrativo nell'utilizzo delle attrezzature informatiche, multimediali e/o audiovisive presenti negli spazi gestiti dall'Università di Udine.
    \par
        La presente Tesi di Laurea, nata dall'elaborazione dell'esperienza maturata in sede di tirocinio curricolare (9 CFU) presso la sede di Ateneo a Gorizia, tratta della figura e del ruolo del Tecnico Informatico all'interno di un Polo Universitario, soffermandosi principalmente sulla collaborazione alla gestione informatica e multimediale nonché sulla manutenzione degli impianti audiovisivi, presenti nel Centro Polifunzionale di Pordenone (CEPO) e nel Centro Polifunzionale di Gorizia (CEGO). \newline
        Per approfondire questo aspetto sono state condotte alcune interviste alle figure-chiave del CEGO che si sono gentilmente rese disponibili all'indagine. 
% ---------------------------------------------------------------------------------------
    \subsection*{Struttura dell'elaborato}
        \par
            L'elaborato proposto è stato organizzato in diversi capitoli, ognuno dei quali è stato suddiviso in paragrafi per meglio fornire le informazioni relative allo sviluppo del tirocinio curricolare e per delineare opportunamente la figura e il ruolo del Tecnico Informatico. 
        \par
            Il primo capitolo, immediatamente successivo a questa introduzione, costituisce la premessa alla descrizione del CEGO e tratta quali siano le figure principali al suo interno attraverso opportune interviste svolte sia in presenza che da remoto tramite la piattaforma Teams. Vengono quindi descritte le differenze tra il CEPO e il CEGO in un paragrafo a sé stante. \newline
            Di seguito sono elencati i dispositivi che si possono incontrare e si descrive il loro ruolo all'interno della topografia fisica degli impianti del CEGO. 
        \par
            Trattando la figura e il ruolo del Tecnico Informatico ne vengono spiegate le varie mansioni, distinguendo tra quelle relative alla collaborazione dalla gestione informatica e multimediale, quelle relative alla manutenzione degli impianti informatici e audiovisivi e, infine, quelle di natura più generale, non facilmente riconducibili a queste due macro categorie.
        \par
            Proseguendo, verrà dedicato un capitolo alla discussione sul futuro della figura del Tecnico Informatico in un mondo digitale sempre più dominato dalle Intelligenze Artificiali. In tale capitolo si illustreranno anche come si potrebbero risolvere le dinamiche tra le due entità, proponendo anche possibili soluzioni di mediate. 
        \par
            Adeguato spazio viene dedicato inoltre ai risultati ottenuti dagli interventi durante il tirocinio e come questi possano impattare la qualità dei servizi offerti all'interno del CEGO, nonché in un Polo Universitario.
            In questa sezione si fa riferimento anche a come il Corso di Laurea in Scienze e Tecnologie Multimediali (CL-20 CL-31) prepari i futuri tecnici e manutentori del settore informatico/multimediale nonché gli insegnanti Tecnico Pratici della classe di concorso B016, Laboratori di scienze e tecnologie informatiche, come previsto dalla normativa vigente.
        \par
            In conclusione alla Tesi di Laurea, vengono riassunti i punti principali della figura del Tecnico Informatico e del suo ruolo all'interno di un Polo Universitario, nonché una riflessione sugli interventi e i dispositivi incontrati.
% ---------------------------------------------------------------------------------------
    \subsection*{Obiettivo}
        \par
            L'obiettivo dell'elaborato è quello di descrivere il ruolo del Tecnico Informatico presso un Polo Universitario e di come il suo ruolo in loco, spesso poco considerato quando si parla di informatica generale, sia invece fondamentale per garantire la continuità dei servizi TIC (Tecnologie dell'Informazione e della Comunicazione) offerti e di garantirne la maggiore qualità possibile a tutti gli utenti. 
% -------------------------------------------------------------------------------------- 
    \subsection*{Contesto e Motivazioni}
        \par
            La figura del Tecnico Informatico è spesso tralasciata nei discorsi di settore in quanto è una figura onnipresente e discreta; tuttavia risulta fondamentale per poter garantire continuità dei servizi offerti e qualità più alta. È altamente necessaria soprattutto in un Polo Universitario, in quanto molti dei servizi che vengono ad oggi offerti risultano composti dalla parte telematica che informatica.
    